% !TeX root = ../Book2.tex
\begin{tikzpicture}[
  concept/.style={draw, rounded corners=2pt, align=center,
    text width=20mm, minimum height=11mm, inner sep=2pt, font=\small},
  every path/.style={semithick}]
\node[concept] (cstar) at (0,0) {\(C^*\)-代数\\与范数};
\node[concept] (spectrum) at (2.6,0) {谱论与\\Gelfand 对偶};
\node[concept] (calculus) at (5.2,0) {连续\\函数演算};
\node[concept] (positive) at (7.8,0) {正性\\与态};
\node[concept] (gns) at (10.4,0) {GNS\\构造};
\node[concept] (faithful) at (13,0) {忠实\\Hilbert 表示};
\draw[->] (cstar) -- (spectrum);
\draw[->] (spectrum) -- (calculus);
\draw[->] (calculus) -- (positive);
\draw[->] (positive) -- (gns);
\draw[->] (gns) -- (faithful);
\node[concept] (operators) at (0,-1.7) {\(B(H)\)};
\node[concept] (topology) at (2.6,-1.7) {SOT 与 WOT};
\node[concept] (bicommutant) at (5.2,-1.7) {双交换子\\定理};
\node[concept] (group) at (7.8,-1.7) {群 von Neumann\\代数};
\node[concept] (extra) at (10.4,-1.7) {纯态、极点\\与近似单位};
\node[concept] (tensor) at (13,-1.7) {张量范数\\极值定理};
\draw[->] (operators) -- (topology);
\draw[->] (topology) -- (bicommutant);
\draw[->] (bicommutant) -- (group);
\draw[->] (faithful) -- node[right,font=\footnotesize] {空间模型} (tensor);
\draw[->] (extra) -- (tensor);
\end{tikzpicture}
